\section{Adaptation to Player-Level Models}
\label{sec:player_level_adaptation}

The models described in the previous sections operate at the club level, where team abilities are modelled directly as team-specific parameters.
However, these models can be adapted to operate at the player level.
This adaptation provides a more granular understanding of team performance by decomposing it into individual player contributions.

Instead of estimating team abilities directly, we estimate individual player abilities and aggregate them, weighting by playing time, to obtain team-level strengths.
For a given match, let $\mathcal{P}_h$ and $\mathcal{P}_a$ denote the sets of players participating in the home and away teams, respectively, and let $m_p$ denote the number of minutes played by player $p$ in that match.
The team-level skill is given by:
\begin{equation}
    \label{eq:team_skill_aggregation}
    s_{\text{team}} = \sum_{p \in \mathcal{P}} \exp(\theta_p) \cdot \frac{m_p}{90},
\end{equation}

\noindent where $\theta_p$ is the log-skill parameter for player $p$, and $\mathcal{P}$ is either $\mathcal{P}_h$ or $\mathcal{P}_a$ depending on the team.
The normalization by 90 minutes ensures that a player's contribution is proportional to their playing time relative to the standard match duration.

When a player is sent off (red card), the team plays with fewer players for the remainder of the match.
To handle this situation, we introduce a synthetic player with an estimable ability parameter $\theta_{\text{synthetic}}$ that represents the impact of playing with fewer players.
Specifically, if a player is sent off at minute $t_{\text{off}}$, we add a synthetic player with $m_{\text{synthetic}} = 90 - t_{\text{off}}$ minutes.
If two or more players receive a red card, the same treatment is applied, increasing the value of $m_{\text{synthetic}}$.
The synthetic player's ability $\theta_{\text{synthetic}}$ is estimated by the model and quantifies the impact of the red card on team performance.
This approach allows the model to learn how much a sending-off affects team strength.

As club-level models, player-level models also require identifiability constraints.
The same sum-to-zero constraint is imposed on the player skill parameters:
\begin{equation}
    \label{eq:player_sum_zero_constraint}
    \sum_{p=1}^{n} \theta_p = 0,
\end{equation}

\noindent where $n$ is the total number of players, including the synthetic.
This constraint ensures that the model is identifiable and that player abilities are measured relative to the average player ability.
